\documentclass[12pt,a4paper]{article}

% ---------- Мова та шрифти ----------
\usepackage{fontspec}
\usepackage{polyglossia}
\setmainlanguage{ukrainian}
\setotherlanguage{english}

\setmainfont{DejaVu Serif}
\setsansfont{DejaVu Sans}
\setmonofont{DejaVu Sans Mono}

% ---------- Компонування ----------
\usepackage[a4paper,margin=2.2cm]{geometry}
\usepackage{graphicx}
\usepackage{float}
\usepackage{caption}
\usepackage{xcolor}
\usepackage{hyperref}
\usepackage{enumitem}
\usepackage{titlesec}
\usepackage{longtable}
\usepackage{array}
\usepackage{booktabs}
\usepackage{microtype}
\setlength{\emergencystretch}{3em}
\sloppy

\hypersetup{
    colorlinks=true,
    linkcolor=blue!50!black,
    urlcolor=blue!50!black,
    citecolor=blue!50!black,
    pdftitle={Порівняння десктоп- та веб-клієнта -- проєкт Remote File Folder Client},
    pdfauthor={Yana Utochkina}
}

\titleformat{\section}{\normalfont\Large\bfseries}{\thesection.}{0.6em}{}
\titleformat{\subsection}{\normalfont\large\bfseries}{\thesubsection}{0.6em}{}

\setlength{\parindent}{0pt}
\setlength{\parskip}{0.6em}

\title{\textbf{Порівняння десктоп- та веб-клієнта\\[2pt] \large проєкту Remote File Folder Client\\(спрощений аналог Google Drive)}}
\author{Уточкіна Яна}
\date{Етап: порівняльний аналіз двох підходів до клієнта \\ \today}

\begin{document}

\maketitle

\begin{abstract}
\noindent
Проєкт \textbf{Remote File Folder Client} постачається з двома клієнтами -- нативним застосунком macOS (Swift/SwiftUI) та браузерним застосунком (React/TypeScript) -- проти одного спільного Go API та бекенду PostgreSQL. Обидва клієнти вже мають власний окремий звіт з описом внутрішньої структури, прецедентів використання та покриття тестами. Цей звіт натомість порівнює їх безпосередньо: що в них спільне, де вони справді відрізняються, і які компроміси несе кожен підхід.
\end{abstract}

\tableofcontents
\newpage

% =====================================================================
\section{Вступ}
% =====================================================================

Обидва клієнти надають ту саму персональну робочу область файлів, звертаються до того самого API через HTTPS за тим самим публічним хостнеймом і відображають ті самі записи файлів. Вони не були створені незалежно одне від одного з подальшим порівнянням постфактум -- функція синхронізації теки веб-клієнта є прямим, свідомим портом рушія синхронізації десктоп-клієнта, адаптованим до того, що реально може браузер. Цей зв'язок формує більшу частину подальшого викладу: обидва клієнти тісно узгоджені щодо того, \emph{що} робить продукт, і відрізняються переважно тим, \emph{як} кожна платформа дозволяє це робити.

% =====================================================================
\section{Спільні елементи}
% =====================================================================

\subsection{Єдиний контракт з бекендом}

Обидва клієнти говорять точно тим самим протоколом з тим самим Go API: реєстрація та вхід за логіном/паролем із JWT access/refresh токенами, та сама структура запису файлу (ім'я, розширення, розмір, часові позначки створення/модифікації, автор вивантаження, редактор), ті самі ендпойнти вивантаження (одноразовий multipart для малих файлів, chunked-сесія initiate/chunk/complete для великих, обидва обмежені, щоб не перевищувати ліміт розміру тіла запиту Cloudflare Tunnel), та той самий ендпойнт попереднього перегляду вмісту \texttt{.java}/\texttt{.png}. Контракт з бекендом жодного з клієнтів не є підмножиною чи надмножиною іншого -- це один і той самий контракт, яким користуються два різні фронтенди.

\subsection{Той самий основний набір функцій}

Реєстрація/вхід, перегляд файлів із сортованими/фільтрованими атрибутами, перегляд \texttt{.java} як тексту та \texttt{.png} як зображення, вивантаження (через вибір файлу \emph{і} шляхом перетягування файлів з робочого столу/браузера), завантаження та видалення -- кожна з цих функцій існує ідентично в обох клієнтах, відрізняючись лише тим, як кожен інструментарій інтерфейсу платформи виражає її (подання SwiftUI проти компонентів React).

\subsection{Той самий алгоритм синхронізації}

Синхронізація теки виконує той самий алгоритм з обох сторін: сканування теки для плоского (нерекурсивного) маніфесту імені/розміру/SHA-256 для кожного файлу, надсилання його разом з останньою синхронізованою версією на \texttt{POST /api/sync/diff}, застосування отриманих дій \texttt{download}/\texttt{delete\_local}/\texttt{conflict}, а потім незалежне вивантаження нових чи змінених локальних файлів і видалення віддалених копій локально видалених. Обидві реалізації виключають ті самі три випадки з цього фінального проходу в межах одного циклу -- файл, уже позначений як конфлікт, файл, який сервер щойно велів видалити, і файл, щойно завантажений у цьому самому циклі -- оскільки пропуск будь-якого з трьох дозволяє проходу неправильно прочитати побічний ефект цього самого циклу як справжню локальну зміну. Обидва розв'язують конфлікт тими самими трьома варіантами (Keep Local, Keep Remote, Keep Both), і обидва називають копію ``keep both'' однаково: \texttt{Name (conflict copy, YYYY-MM-DD).ext}.

\subsection{Збережений курсор синхронізації}

Обидва клієнти пам'ятають номер останньої синхронізованої версії та запис по кожному файлу (віддалений ID, розмір, хеш) між перезапусками, тож синхронізації не доводиться щоразу починати з нуля. Для цього використовуються різні технології зберігання (див. \S\ref{sec:differences}), але дані, які зберігаються, і мета їх зберігання ідентичні.

% =====================================================================
\section{У чому вони відрізняються}
\label{sec:differences}
% =====================================================================

\subsection{Тригер синхронізації: безперервний проти разового}

Десктоп-клієнт спостерігає за обраною текою у фоновому режимі через \texttt{FSEvents} і додатково опитує сервер кожні 45 секунд, тож зміна з будь-якого боку помічається за мить без будь-яких дій користувача. Браузер узагалі не може безперервно спостерігати за локальною текою у фоновому режимі -- тож синхронізація веб-клієнта є одним разовим кліком кнопки \textbf{Sync now}: один цикл узгодження за запитом, і між кліками нічого не відбувається, незалежно від того, що змінюється з будь-якого боку.

\subsection{Доступ до теки: нативні шляхи проти піскового дескриптора}

Десктоп-клієнт працює зі звичайними шляхами файлової системи через \texttt{FileManager} -- після обрання теки її можна вільно й безстроково читати й записувати, без подальших запитів дозволу. Веб-клієнт може дістатися лише теки, яку користувач явно надав через File System Access API браузера, який повертає непрозорий, пісочний дескриптор, а не шлях; дозвіл для цього дескриптора може втратити чинність між сесіями і може бути повторно запитаний лише зсередини справжнього кліку користувача, а не автоматично при завантаженні сторінки. Сам API доступний лише в браузерах на основі Chromium (Chrome, Edge, Brave) -- у Firefox та Safari немає жодного еквівалента, тож уся функція синхронізації там недоступна, тоді як решта функцій (перегляд, вивантаження, завантаження, видалення) працює в будь-якому браузері.

\subsection{Сортування й фільтрація: на клієнті проти на сервері}

Десктоп-клієнт завантажує повний список файлів користувача один раз (до фіксованої межі) і сортує/фільтрує його локально в пам'яті при кожній взаємодії з інтерфейсом. Веб-клієнт натомість надсилає активний стовпець сортування, порядок і фільтр розширення на сервер як параметри запиту й повторно завантажує пагіновану сторінку при кожній зміні, ніколи не тримаючи в пам'яті більш ніж віконний (віртуалізований) набір рядків одночасно. Обидва дають той самий видимий результат -- вони просто розміщують роботу сортування/фільтрації на протилежних кінцях з'єднання.

\subsection{Зберігання локального стану}

Десктоп-клієнт тримає курсор синхронізації та відомі записи файлів в одному JSON-файлі в Application Support, який читається й записується через \texttt{Codable}. Веб-клієнт тримає такі самі дані в IndexedDB, за невеликою абстракцією ключ-значення, оскільки сторінка браузера не має власної файлової системи, куди можна записати звичайний файл -- і там само зберігає дескриптор обраної теки, що взагалі можливо лише тому, що \texttt{FileSystemDirectoryHandle} придатний до структурного клонування.

\subsection{Мережевий стек}

Десктоп-клієнт використовує \texttt{URLSession} для кожного запиту, включно з прогресом chunked-вивантаження. Веб-клієнт використовує \texttt{fetch} для всього, крім звітування про прогрес chunked-вивантаження, яке виконує через \texttt{XMLHttpRequest}, оскільки \texttt{fetch} взагалі не має події прогресу вивантаження.

\subsection{Поширення та охоплення}

Десктоп-клієнт постачається як \texttt{.dmg}, потребує одноразового встановлення, працює лише на macOS і потребує повторного встановлення для отримання нової версії. Веб-клієнт не потребує жодного встановлення, працює на будь-якому пристрої з сучасним браузером і завжди є найновішою версією в момент завантаження сторінки -- за ціну того, що саме функція синхронізації обмежена браузерами на основі Chromium, як описано вище.

\subsection{Модель конкурентності}

Оскільки синхронізацію десктоп-клієнта можуть одночасно запустити два незалежних джерела -- опитування кожні 45 секунд та відкладений callback спостереження за текою -- його процедура узгодження має захищатися від одночасного повторного запуску; без цього захисту два перетинних проходи, що читають і записують той самий стан, можуть кожен працювати зі застарілим знімком змін, що вже виконує інший. Веб-клієнт не має такого ризику за своєю конструкцією: прохід синхронізації завжди починається лише з одного кліку кнопки, тож немає другого тригера, з яким можна б було конкурувати.

% =====================================================================
\section{Компроміси кожного підходу}
% =====================================================================

\subsection{Десктоп (Swift/SwiftUI, нативний macOS)}

\textbf{Переваги:}
\begin{itemize}[nosep]
    \item Безперервна фонова синхронізація -- зміни з будь-якого боку з'являються без жодних дій користувача;
    \item Необмежений, бездозвільний доступ до будь-якої локальної теки після її обрання;
    \item Нативна інтеграція з ОС: Keychain для зберігання токенів, нативні діалоги вибору файлів, без пісочниці браузера, яку потрібно обходити;
    \item Жодних винятків через сумісність браузерів -- кожна функція працює однаково щоразу.
\end{itemize}

\textbf{Недоліки:}
\begin{itemize}[nosep]
    \item Лише для macOS, і потребує встановлення (та періодичного повторного встановлення) \texttt{.dmg};
    \item Сортування/фільтрація тримає весь список файлів користувача в пам'яті й обробляє його на клієнті, що нормально в масштабі персонального диска, але не масштабується безмежно;
    \item Модель безперервної синхронізації архітектурно складніша, і ця складність уже спричинила три реальні помилки конкурентності/застарілості під час розробки (витік стану при перемиканні теки, перегони через повторний вхід у процедуру узгодження та видалення файлу в тому самому циклі, що й завантаження), які довелося знайти та виправити.
\end{itemize}

\subsection{Веб (React/TypeScript, браузер)}

\textbf{Переваги:}
\begin{itemize}[nosep]
    \item Без встановлення, працює на будь-якій ОС із сучасним браузером, завжди найновіша версія;
    \item Серверне сортування/фільтрація та пагінація масштабуються до значно більшої кількості файлів, ніколи не тримаючи весь список на клієнті;
    \item Суттєво простіша модель конкурентності синхронізації -- одна кнопка, один прохід, немає фонового тригера, з яким можна конкурувати;
    \item Успадкував уже налагоджений алгоритм синхронізації, тож запустився з усіма трьома важко здобутими винятками застарілості десктоп-клієнта вже на місці, а не з необхідністю повторно виявляти їх самостійно.
\end{itemize}

\textbf{Недоліки:}
\begin{itemize}[nosep]
    \item Жодної фонової синхронізації -- зміна, зроблена деінде, локально невидима, доки користувач не згадає повторно клікнути \textbf{Sync now};
    \item Саме функція синхронізації працює лише в браузерах на основі Chromium; вона повністю відсутня у Firefox та Safari;
    \item Дозвіл на теку не постійний -- він може втратити чинність між сесіями і має бути наданий знову зі свіжого кліку користувача, що додає тертя, якого немає в нативному застосунку;
    \item Збережений стан живе в сховищі браузера (IndexedDB), яке підпорядковується власній політиці витіснення сховища браузера під тиском дискового простору, на відміну від власних файлів нативного застосунку на диску.
\end{itemize}

% =====================================================================
\section{Підсумкова таблиця}
% =====================================================================

Контракт з бекендом, основний набір функцій та алгоритм синхронізації/модель конфліктів ідентичні для обох клієнтів (див.\ розд.~2), тож не включені в таблицю нижче -- таблиця~\ref{tab:summary} стисло зводить лише ті аспекти, де обидва клієнти справді відрізняються.

\small
\renewcommand{\arraystretch}{1.3}
\begin{longtable}{@{}p{2.7cm}p{6.35cm}p{6.35cm}@{}}
\caption{Десктоп проти веб-клієнта, одним поглядом}
\label{tab:summary} \\
\toprule
\textbf{Аспект} & \textbf{Десктоп (Swift/SwiftUI)} & \textbf{Веб (React/TypeScript)} \\
\midrule
\endhead
Платформа & Лише macOS & Будь-яка ОС із сучасним браузером \\
Тригер синхронізації & Безперервний -- \texttt{FSEvents} + опитування кожні 45с & Разовий -- один клік, один прохід \\
Доступ до теки & Нативний шлях, бездозвільний після обрання & Пісочний дескриптор; дозвіл може втратити чинність \\
Доступність синхронізації & Завжди (нативний застосунок) & Лише Chromium -- відсутня у Firefox/Safari \\
Сортування й фільтрація & На клієнті, над обмеженим повним завантаженням & На сервері, через параметри пагінованого запиту \\
Зберігання стану & JSON-файл & IndexedDB, за абстракцією ключ-значення \\
Мережевий стек & \texttt{URLSession} & \texttt{fetch} + \texttt{XMLHttpRequest} \\
Поширення & Встановлюваний \texttt{.dmg}; повторне встановлення & Без встановлення; завжди найновіша версія \\
Ризик конкурентності & Потребує явного захисту від повторного входу & Відсутній за конструкцією \\
\bottomrule
\end{longtable}
\renewcommand{\arraystretch}{1}

% =====================================================================
\section{Висновок}
% =====================================================================

Жоден з клієнтів не є просто кращою версією іншого -- вони обмінюють те саме обмеження (браузер не може спостерігати за текою у фоновому режимі) на два різні набори переваг. Десктоп-клієнт -- це місце, де справді доречна безперервна, безтертя синхронізація теки, за ціну нативного, встановлюваного застосунку лише для macOS. Веб-клієнт -- це місце, де доречний доступ без встановлення, кросплатформний і завжди найновіший, за ціну того, що синхронізація є разовою дією, обмеженою специфічним для Chromium API браузера. Велика область перетину -- той самий контракт з бекендом, той самий основний набір функцій, а саме для синхронізації -- точно той самий алгоритм і модель конфліктів -- це те, що робить обидва схожими на один продукт із двома вхідними дверима, а не на два непов'язані застосунки, які просто випадково користуються спільним сервером.

\end{document}
